% PREÂMBULO - CONFIGURAÇÕES ABNT (CLASSE ARTICLE)
% ====================================================
% \usepackage{abntex2}
\usepackage[utf8]{inputenc}
\usepackage[T1]{fontenc}
% \usepackage[brazilian]{babel}
\usepackage{lipsum}
\usepackage{times}
\usepackage{geometry}
\usepackage{setspace}
\usepackage[alf,abnt-emphasize=bf]{abntex2cite}
\usepackage{titlesec}
\usepackage{titletoc}
\usepackage{caption}
\usepackage{etoolbox}
\usepackage{amsmath}
\usepackage{changepage}
\usepackage{ragged2e}
\usepackage{siunitx}

\setlength{\headheight}{14.5pt}

% --- PACOTES GRÁFICOS ESSENCIAIS ---
\usepackage{graphicx} % Figuras básicas
\usepackage{float}    % Posicionamento preciso
\usepackage{booktabs} % Tabelas profissionais
\usepackage{subcaption} % Subfiguras
\usepackage{caption} % Subfiguras
\usepackage{pgfplots}  % Gráficos avançados
\pgfplotsset{compat=1.18}
\usepackage{tikz}      % Diagramas vetoriais
\usepackage{multicol}

% --- MARGENS ABNT ---
\geometry{a4paper,left=3cm,right=2cm,top=3cm,bottom=2cm}

\onehalfspacing                 % Espaçamento 1,5 linha (ABNT)

% --- RECUO E ESPAÇAMENTO PADRÃO ABNT ---
\setlength{\parindent}{1.25cm} % Recuo de parágrafo
\setlength{\parskip}{0pt} % Remove espaço extra entre parágrafos
\raggedbottom % Evita espaçamento irregular


\usepackage{fancyhdr}
\fancyhf{} % limpa todos os campos do cabeçalho e rodapé
\fancyhead[R]{\thepage} % número da página à direita no topo
\renewcommand{\headrulewidth}{0pt} % remove linha horizontal do topo
\renewcommand{\footrulewidth}{0pt} % remove linha no rodapé

% Configuração da fonte
\renewcommand{\paragraph}{\fontsize{12}{14}\selectfont} % Tamanho 12 com espaçamento 1.5

% --- CORREÇÃO DO RECUO NO PRIMEIRO PARÁGRAFO ---
\usepackage{indentfirst} % Garante recuo mesmo no primeiro parágrafo
\setlength{\parindent}{1.25cm} % Recuo ABNT padrão

% --- FORMATAÇÃO DAS SEÇÕES (ABNT) ---
\usepackage{titlesec}

% Seções (Nível 1)
\titleformat{\section}
  [hang] % alinhado à esquerda
  {\normalfont\bfseries\fontsize{12}{14}\selectfont\MakeUppercase}
  {\thesection}{1em}{}


% Subseções (Nível 2)
\titleformat{\subsection}
  [hang]
  {\normalfont\fontsize{12}{14}\selectfont}
  {\thesubsection}{1em}{}


% Subsubseções (Nível 3)
\titleformat{\subsubsection}
  [hang]
  {\normalfont\fontsize{12}{14}\selectfont}
  {\thesubsubsection}{1em}{}

% --- COMANDO PERSONALIZADO PARA "CAPÍTULOS" NÃO NUMERADOS ---
\newcommand{\capitulo}[1]{
  \begin{center}
    \vspace{2em}
    \textbf{\MakeUppercase{#1}}
    \vspace{1em}
  \end{center}
}


% Espaçamento vertical
\titlespacing*{\section}{0pt}{18pt}{6pt}    % [left]{before}{after}
\titlespacing*{\subsection}{0pt}{12pt}{6pt}
\titlespacing*{\subsubsection}{0pt}{12pt}{6pt}

% Profundidade de numeração
\setcounter{secnumdepth}{3} % Até subsubseções

\usepackage[titles]{tocloft}

% --- FORMATAÇÃO DO SUMÁRIO ---
% \usepackage{tocloft}
% \usepackage[nottoc]{tocbibind} % Evita que o título "Sumário" apareça no próprio sumário
\renewcommand{\contentsname}{\begin{center}
        \bfseries\MakeUppercase{Sumário}
      \end{center}}
% Recuo das seções (ABNT padrão: sem recuo)
\titlecontents{section}[1.5em]           % Recuo da seção
{\addvspace{0pt}\bfseries\normalsize}  % Estilo
{\contentslabel{1.5em}}                % Alinhamento do número
{}                                     % Sem recuo adicional
{\titlerule*[0.6pc]{.}\contentspage}   % Pontilhado e número de página


% Recuo das subseções (ligeiramente menor)
\titlecontents{subsection}[1.5em]
{\normalfont\normalsize}
{\contentslabel{1.5em}}
{}
{\titlerule*[0.6pc]{.}\contentspage}



% Centraliza e coloca em maiúsculas a lista de figuras
\renewcommand{\listfigurename}{
  \bfseries\MakeUppercase{Figuras}}

% Centraliza e coloca em maiúsculas a lista de tabelas
\renewcommand{\listtablename}{\begin{center}
  \bfseries\MakeUppercase{Lista de Tabelas}
\end{center}}

% --- CONFIGURAÇÕES DE FIGURAS E TABELAS ---
\captionsetup[figure]{
  position=above,    % Legenda acima da figura
  justification=centering,
  labelsep=endash,
  name=Figura,
  font=small,
  labelfont=bf
}

% Comando personalizado para fonte abaixo
\newcommand{\fonte}[1]{
  \vspace{-0.5em} % Ajuste de espaçamento
  \begin{center}
    \footnotesize\textbf{Fonte:} #1
  \end{center}
}

% --- CONFIGURAÇÃO DE TABELAS ABNT ---
\captionsetup[table]{
  position=above,     % Legenda acima da tabela
  justification=centering,
  labelsep=endash,
  name=Tabela,
  font=small,
  labelfont=bf
}

\newcommand{\fontet}[1]{
  \begin{center}
    \footnotesize\textbf{Fonte:} #1
  \end{center}
}

% Comandos personalizados para o formato ABNT
\DeclareCaptionLabelFormat{abntformat}{\textbf{#1~#2}}
\captionsetup[figure]{labelformat=abntformat}
\captionsetup[table]{labelformat=abntformat}

\usepackage{tocloft}

% --- CONFIGURAÇÕES PARA LISTA DE FIGURAS ---
\setlength{\cftfigindent}{0pt}
\setlength{\cftfignumwidth}{6em}
\renewcommand{\cftfigpresnum}{\textbf{FIGURA~}}
\renewcommand{\cftfigaftersnum}{\textbf{~--~}}
\renewcommand{\cftfigdotsep}{3}

% --- CONFIGURAÇÕES PARA LISTA DE TABELAS ---
\setlength{\cfttabindent}{0pt}             % Sem recuo
\setlength{\cfttabnumwidth}{6em}           % Espaço reservado para "TABELA 10"
\renewcommand{\cfttabpresnum}{\textbf{TABELA~}}  % Prefixo
\renewcommand{\cfttabaftersnum}{\textbf{~--~}}   % Sufixo
\renewcommand{\cfttabdotsep}{3}            % Mais pontos entre nome e número de página



% Define novo tipo de float
% \newfloat{myequation}{htbp}{loe}
% \floatname{myequation}{EQUAÇÃO}

% % Configura legenda do float
% \captionsetup[myequation]{
%   labelfont=bf,
%   labelsep=space,
%   justification=centering,
%   singlelinecheck=false
% }

% Ambiente "simbolos" no estilo ABNT
\usepackage{enumitem}

\newlist{simbolos}{description}{1}
\setlist[simbolos]{%
  labelwidth=1cm,        % Largura da coluna do símbolo
  labelsep=1cm,          % Espaço entre símbolo e descrição
  leftmargin=2cm,        % Recuo à esquerda (soma de labelwidth + labelsep)
  style=nextline,        % Quebra de linha após o símbolo
  font=\normalfont       % Fonte normal no símbolo
}

% Cria a lista de gráficos
\newlistof{grafico}{lol}{Gráficos}

% Formatação da numeração da legenda no corpo do documento
\renewcommand{\thegrafico}{\arabic{grafico}}

% Formatação da numeração na lista de gráficos (prefixo em maiúsculas e negrito)
\renewcommand{\cftgraficopresnum}{\textbf{GRÁFICO~}}
\renewcommand{\cftgraficoaftersnum}{\textbf{~--~}}  % traço com espaços antes e depois
\setlength{\cftgraficonumwidth}{7em}      % largura reservada para o número
\setlength{\cftgraficoindent}{0pt}        % sem indentação extra

% Configura legenda para o tipo 'grafico'
\captionsetup[grafico]{
  position=above,
  justification=centering,
  labelsep=endash,
  font=small,
  labelfont=bf,
  name=Gráfico,
  labelformat=simple
}

\setcounter{grafico}{0}

% Comando para inserir gráfico com legenda acima, fonte abaixo e label
\newcommand{\inserirgrafico}[5][]{%
  \begin{center}
    \captionsetup{type=grafico,position=above}%
    \captionof{grafico}{#3}%
    \label{#5}% label logo após a legenda
    \includegraphics[#1]{#2}%
    \par\vspace{0.1em}%
    {\footnotesize\textbf{Fonte:} #4}%
  \end{center}
}

% Define o nome da lista de equações
\newcommand{\listequationsname}{Equações}

% Cria a lista de equações com extensão .loe
\newlistof{equationlist}{loe}{\listequationsname}

% Prefixo e sufixo na lista de equações (prefixo em negrito e traço)
\renewcommand{\cftequationlistpresnum}{\textbf{EQUAÇÃO~}}
\renewcommand{\cftequationlistaftersnum}{\textbf{~--~}}
\setlength{\cftequationlistnumwidth}{7em}
\setlength{\cftequationlistindent}{0pt}

% Comando para adicionar entrada na lista de equações
% #1 = texto da entrada na lista (normalmente a equação ou descrição)
\newcommand{\addtoequationlist}[1]{%
  \addcontentsline{loe}{equationlist}{\protect\numberline{\theequation}#1}%
}





% --- Ambiente Personalizado ErrataABNT ---
% ==== NO PREÂMBULO ====  
\usepackage{geometry}   % para controle de margem  
\usepackage{setspace}   % espaçamento  
\usepackage{array}      % para tabelas  
\usepackage{caption}    % legendas  

% Comando para inserir a página de Errata ABNT
\newcommand{\paginaErrata}[1]{%
  \clearpage
  \thispagestyle{empty}         % sem cabeçalho/rodapé  
  \newgeometry{left=1cm, right=2cm, top=3cm, bottom=2cm}  
  \section*{\begin{center}Errata\end{center}}        % título senúmero  
  % \vspace{1\baselineskip}  
  \setstretch{1.0}               % espaçamento simples  
  % ==== TABELA DE ERROS ====  
  {\renewcommand{\arraystretch}{1.5}% altura das linhas
  \input{tabelas/erros}
  % \begin{tabular}{@{} p{2.5cm} p{3cm} p{5cm} p{5cm} @{}}
  %   \hline
  %   \textbf{Folha} & \textbf{Linha} & \textbf{Onde se lê} & \textbf{Leia-se} \\
  %   \hline
  %   10             & 15             & \textit{“esperem resultados imediatos”} & \textit{“esperem resultados mediatos”} \\
  %   45             & 3              & “processo contínuo de ensino e apredizado” & “processo contínuo de ensino e aprendizado” \\
  %   % adicione quantas linhas forem necessárias
  %   \hline
  % \end{tabular}
  }
  \restoregeometry
  \clearpage
}


% --- CONFIGURAÇÕES PARA O SUMÁRIO ---
\setlength{\cftsecindent}{0pt}             % Seções alinhadas à esquerda
\setlength{\cftsecnumwidth}{0em}
\renewcommand{\cftsecleader}{\cftdotfill{\cftdotsep}} % Aplica pontilhado
\renewcommand{\cftdotsep}{3}               % Aumenta pontos do sumário

% (Opcional) para subseções também
\setlength{\cftsubsecindent}{0pt}
\setlength{\cftsubsecnumwidth}{4em}
\renewcommand{\cftsubsecleader}{\cftdotfill{\cftdotsep}}
\renewcommand{\cftdotsep}{3}               % Aumenta pontos do sumário

% Pré-requisitos
\usepackage{graphicx}  % rotatebox
\usepackage{geometry}  % definir dimensões da página

\newcommand{\lombada}[3]{%
  \clearpage
  \newgeometry{left=0cm, top=0cm, bottom=0cm, right=0cm}
  \thispagestyle{empty}

  % Ambiente centralizado vertical com largura total da lombada
  \begin{center}
    \vspace*{1cm} % Espaço até autor (posição vertical)
    \rotatebox{270}{{\textbf{\MakeUppercase{#1}}}} % Autor
    
    \vspace*{2cm} % Espaço até o título (aproximadamente até 15cm do topo)
    \rotatebox{270}{{\textbf{\MakeUppercase{#2}}}} % Título

      % Último item: horizontal, alinhado à esquerda, com largura de 4cm
  \vspace*{2.5cm} % Espaço até parte inferior (29.7 - 26.5 ≈ 3.2)
  \hspace*{0.4cm} % margem esquerda interna da lombada (ajustável)
  \parbox[t]{5cm}{\centering\normalsize\MakeUppercase{#3}}

  \end{center}
  
  \restoregeometry
  \clearpage
}

% Ficha catalográfica formatada segundo a ABNT
\usepackage{geometry} % já deve estar incluso

\newcommand{\fichacatalografica}[9]{%
  \clearpage
  \newgeometry{left=3cm, right=2cm, top=3cm, bottom=2cm}
  \thispagestyle{empty}
  \vspace*{\fill}
  \hspace*{\fill}
  \fbox{%
    \begin{minipage}[t][7.5cm][t]{12.5cm}
      \setstretch{1.0}
      \normalfont\fontsize{12}{14}\selectfont
      \raggedright
      #1\\ % Autor
      \MakeUppercase{#2} / #1. – #3 ed. – #4: #5, #6.\\ p. 35\\
      \textbf{Classificação:} #7\\
      \textbf{ISBN:} #8
      \textbf{CDD:} #9
    \end{minipage}
  }%
  \vspace*{1cm}
  \restoregeometry
  \clearpage
}





% --- FORMATAÇÃO DA BIBLIOGRAFIA ABNT COM TÍTULO CENTRALIZADO, NEGRITO E UPPERCASE ---
\makeatletter
\renewenvironment{thebibliography}[1]
     {
      \begin{center}
        \bfseries\MakeUppercase{Referências}
      \end{center}
      \@mkboth{\MakeUppercase{\refname}}{\MakeUppercase{\refname}}%
      \list{\@biblabel{\@arabic\c@enumiv}}%
           {\leftmargin=0cm
            \labelwidth=0cm
            \labelsep=0cm
            \itemindent=0pt
            \parsep=0pt
            \itemsep=1.5\baselineskip
            \usecounter{enumiv}}%
      \singlespacing
      \sloppy
      \clubpenalty4000
      \@clubpenalty \clubpenalty
      \widowpenalty4000%
      \sfcode`\.\@m
     }
     {\def\@noitemerr
       {\@latex@warning{Empty `thebibliography' environment}}%
      \endlist}
\makeatother